\documentclass[12pt, letterpaper, oneside]{report}

% ------------------------------------------------------------
% PAQUETES
% ------------------------------------------------------------
\usepackage[utf8]{inputenc}
\usepackage[spanish, es-tabla]{babel}
\usepackage{mathptmx}
\usepackage[letterpaper, left=4cm, top=3cm, right=3cm, bottom=3cm]{geometry}
\usepackage{setspace}
\usepackage{indentfirst}
\usepackage{fancyhdr}
\usepackage{titlesec}
\usepackage{tikz}
\usepackage{graphicx}
\usepackage{tocloft}
\usepackage[bookmarks=false, colorlinks=true, linkcolor=black, urlcolor=black, citecolor=black]{hyperref}
\usetikzlibrary{shapes, positioning}

% ------------------------------------------------------------
% FORMATO GENERAL
% ------------------------------------------------------------
\setstretch{1.5}
\setlength{\parindent}{0.5cm}
\setlength{\parskip}{0pt}

% ------------------------------------------------------------
% FORMATO DE CAPÍTULOS
% ------------------------------------------------------------
\titleformat{\chapter}[display]
  {\centering\normalfont\bfseries\large}
  {CAPÍTULO \Roman{chapter}}
  {0.5em}
  {}
\titlespacing*{\chapter}{0pt}{-20pt}{20pt}
\titlespacing*{\section}{0pt}{12pt}{6pt}
\titlespacing*{\subsection}{0pt}{6pt}{3pt}

% ------------------------------------------------------------
% ENCABEZADO Y PIE DE PÁGINA
% ------------------------------------------------------------
\fancypagestyle{preliminary}{%
  \fancyhf{}
  \fancyfoot[C]{\thepage}
  \renewcommand{\headrulewidth}{0pt}
}
\fancypagestyle{body}{%
  \fancyhf{}
  \fancyhead[R]{\thepage}
  \renewcommand{\headrulewidth}{0pt}
}

% ------------------------------------------------------------
% TABLA DE CONTENIDOS
% ------------------------------------------------------------
\renewcommand{\cfttoctitlefont}{\hfill\bfseries\large}
\renewcommand{\cftaftertoctitle}{\hfill}
\setlength{\cftbeforetoctitleskip}{-20pt}
\renewcommand{\cftchapfont}{\normalfont}
\renewcommand{\cftchappagefont}{\normalfont}
\setlength{\cftchapindent}{0pt}
\setlength{\cftsecindent}{0.5cm}
\setlength{\cftsubsecindent}{1cm}

% ------------------------------------------------------------
% DOCUMENTO
% ------------------------------------------------------------
\begin{document}

\pagenumbering{roman}

% ============================================================
% PORTADA
% ============================================================
\begin{titlepage}
\thispagestyle{empty}
\noindent\makebox[\textwidth][c]{\llap{\raisebox{-0.6cm}{\includegraphics[width=2.8cm]{../logoIUTECP_azul.png}}\hspace{0.2cm}}%
\begin{tabular}[t]{c}
  \small
  REP\'UBLICA BOLIVARIANA DE VENEZUELA\\
  MINISTERIO DEL PODER POPULAR PARA\\
  LA EDUCACI\'ON UNIVERSITARIA\\
  INSTITUTO UNIVERSITARIO DE TECNOLOG\'IA\\
  ``EL\'IAS CALIXTO POMPA''\\
  INFORM\'ATICA\\
  AN\'ALISIS Y DISE\~NO DE SISTEMAS II
\end{tabular}}

\vfill

\begin{center}
  \large\textbf
  {SISTEMA DE CONTROL DE ASISTENCIA\\MEDIANTE C\'ODIGOS QR PARA EL IUTECP\\[0.6cm]
  Casos de Uso del Sistema}
\end{center}

\vfill

\begin{flushright}
  \large
  \textbf{Facilitador:} William Hern\'andez

  \vspace{0.3cm}

  \textbf{Integrantes:}\\[0.4cm]
  \begin{tabular}{l l}
    Fernando Rodr\'iguez & C.I: 31.853.086\\
    Daniel Dias          & C.I: 30.937.908
  \end{tabular}
\end{flushright}

\vspace{0.5cm}

\begin{center}
  Guatire, 14 de Agosto de 2026
\end{center}
\end{titlepage}

% ============================================================
% PÁGINAS PRELIMINARES
% ============================================================
\pagestyle{preliminary}

\tableofcontents
\newpage

% ------------------------------------------------------------
% RESUMEN
% ------------------------------------------------------------
\chapter*{RESUMEN}
\addcontentsline{toc}{chapter}{RESUMEN}

El presente informe describe los casos de uso del Sistema de Control de Asistencia mediante C\'odigos QR para el Instituto Universitario de Tecnolog\'ia ``El\'ias Calixto Pompa'' (IUTECP). Se identifican y describen los actores que interact\'uan con el sistema ---administrador, profesor y estudiante---, as\'i como el conjunto de casos de uso que representan las funcionalidades principales, entre las que destacan el inicio de sesi\'on, la generaci\'on y el escaneo de c\'odigos QR, la gesti\'on de usuarios y la consulta de reportes de asistencia. El documento incluye el modelo gr\'afico de casos de uso y la especificaci\'on detallada de cada caso con sus precondiciones, postcondiciones y flujos de eventos.

\newpage

% ============================================================
% CUERPO DEL TRABAJO
% ============================================================
\pagenumbering{arabic}
\pagestyle{body}

% ============================================================
% CAPÍTULO I
% ============================================================
\chapter{DESCRIPCI\'ON DE LA FUNCIONALIDAD DEL SISTEMA}

El control de asistencia en instituciones educativas constituye un proceso fundamental para la gesti\'on acad\'emica. Tradicionalmente, este proceso se realiza mediante listas impresas que el profesor debe portar y marcar manualmente durante cada clase, lo que consume tiempo valioso y est\'a sujeto a errores humanos, suplantaci\'on de identidad y p\'erdida de registros. En el contexto del Instituto Universitario de Tecnolog\'ia ``El\'ias Calixto Pompa'' (IUTECP), se hace necesario contar con una soluci\'on tecnol\'ogica que optimice este proceso.

\section{Prop\'osito del Sistema}

El sistema propuesto tiene como objetivo automatizar el registro de asistencia de los estudiantes mediante el uso de c\'odigos QR generados din\'amicamente por el profesor durante cada sesi\'on de clase. La soluci\'on permitir\'a:

\begin{itemize}
  \item Agilizar el proceso de toma de asistencia, reduci\'endolo a segundos.
  \item Proveer reportes en tiempo real sobre el ausentismo estudiantil.
  \item Centralizar la informaci\'on en una base de datos \'unica accesible desde cualquier dispositivo.
\end{itemize}

\section{Arquitectura General}

El sistema se concibe bajo una arquitectura cliente-servidor utilizando tecnolog\'ias web est\'andar. El backend se desarrolla en PHP 8.4 con conexiones seguras a una base de datos MariaDB mediante sentencias preparadas PDO, mientras que el frontend emplea HTML5, CSS3 y JavaScript vanilla para garantizar compatibilidad y rendimiento sin dependencias externas.

\section{Funcionalidades por Rol}

\subsection{Administrador}

El administrador contar\'a con un panel integral de gesti\'on que incluye:

\begin{itemize}
  \item Gesti\'on completa de usuarios: creaci\'on, activaci\'on y desactivaci\'on de cuentas de profesores y estudiantes.
  \item Administraci\'on de cursos: asignaci\'on de profesores a secciones, creaci\'on y edici\'on de cursos.
  \item Gesti\'on de periodos acad\'emicos y carga masiva de inscripciones desde archivos CSV.
  \item Consulta de reportes de asistencia filtrados por curso, semana o estudiante.
  \item Monitoreo de ausentismo con alertas autom\'aticas para estudiantes en riesgo (25\,\% a 30\,\% de inasistencias) y cr\'itico (m\'as de 30\,\%).
  \item Visualizaci\'on de registros de auditor\'ia con los \'ultimos eventos del sistema.
  \item Generaci\'on de respaldos de la base de datos.
\end{itemize}

\subsection{Profesor}

El profesor dispondr\'a de las siguientes capacidades:

\begin{itemize}
  \item Visualizaci\'on de los cursos que tiene asignados en el periodo activo.
  \item Generaci\'on de c\'odigos QR \'unicos para cada sesi\'on de clase, con expiraci\'on configurable entre 1 y 60 minutos.
  \item Rotaci\'on autom\'atica del c\'odigo QR cada cierto intervalo de segundos (de 10 a 300), lo que impide la reutilizaci\'on del c\'odigo por parte de terceros.
  \item Consulta en tiempo real del estado de asistencia de sus estudiantes.
\end{itemize}

\subsection{Estudiante}

El estudiante podr\'a:

\begin{itemize}
  \item Visualizar las materias en las que se encuentra inscrito junto con su progreso en el pensum de la carrera.
  \item Escanear el c\'odigo QR proyectado por el profesor utilizando la c\'amara de su dispositivo m\'ovil.
  \item Registrar su asistencia de forma autom\'atica, siempre que el c\'odigo QR se encuentre vigente.
  \item Consultar su historial de asistencia.
\end{itemize}

\newpage

% ============================================================
% CAPÍTULO II
% ============================================================
\chapter{ACTORES DEL SISTEMA}

Los actores representan los roles que interact\'uan externamente con el sistema. Cada actor agrupa un conjunto de responsabilidades y funcionalidades a las que puede acceder seg\'un su rol.

\section{Identificaci\'on de Actores}

Se han identificado los siguientes actores:

\begin{description}
  \item[Administrador:] Usuario con privilegios totales sobre el sistema. Se encarga de la gesti\'on de usuarios, cursos, periodos acad\'emicos y la supervisi\'on general del sistema, as\'i como de la consulta de reportes de asistencia.
  \item[Profesor:] Usuario responsable de generar los c\'odigos QR durante sus sesiones de clase, configurar los par\'ametros de asistencia y consultar el estado de asistencia de sus estudiantes.
  \item[Estudiante:] Usuario que registra su asistencia mediante el escaneo de c\'odigos QR y consulta su historial de asistencia.
  \item[Sistema:] Actor secundario que representa los procesos autom\'aticos del sistema, como la validaci\'on de expiraci\'on del QR, la generaci\'on de reportes y el registro de auditor\'ia.
\end{description}

\section{Matriz de Responsabilidades por Actor}

La siguiente tabla resume las responsabilidades principales de cada actor del sistema:

\begin{center}
\begin{tabular}{|p{3cm}|p{10.4cm}|}
\hline
\textbf{Actor} & \textbf{Responsabilidades} \\
\hline
Administrador & Gesti\'on de usuarios, cursos y periodos; monitoreo de ausentismo; consulta de reportes; gesti\'on de respaldos y auditor\'ia \\
\hline
Profesor & Generaci\'on de c\'odigos QR; configuraci\'on de expiraci\'on y rotaci\'on; consulta de asistencias en tiempo real \\
\hline
Estudiante & Escaneo de c\'odigos QR; registro de asistencia; consulta de historial y materias inscritas \\
\hline
Sistema & Validaci\'on de QR y expiraci\'on; generaci\'on de reportes; registro de auditor\'ia; control de accesos \\
\hline
\end{tabular}
\end{center}

\newpage

% ============================================================
% CAPÍTULO III
% ============================================================
\chapter{MODELO DE CASOS DE USO}

El modelo de casos de uso permite representar gr\'aficamente las interacciones entre los actores externos y las funcionalidades que el sistema debe ofrecer.

\section{Diagrama de Casos de Uso}

A continuaci\'on se presenta la representaci\'on gr\'afica del modelo de casos de uso, elaborada con la notaci\'on UML est\'andar:

\begin{center}
\setlength{\fboxsep}{12pt}
\begin{tikzpicture}[
    usecase/.style={ellipse, draw, align=center, minimum width=3cm, minimum height=0.8cm, font=\small},
    every node/.style={font=\small}
  ]

  % ---- ACTORES (STICK FIGURES) ----
  \coordinate (adminPos) at (0,0);
  \draw[thick] (adminPos) circle (0.2);
  \draw[thick] (0,-0.2) -- (0,-0.7);
  \draw[thick] (-0.25,-0.4) -- (0.25,-0.4);
  \draw[thick] (0,-0.7) -- (-0.2,-1.1);
  \draw[thick] (0,-0.7) -- (0.2,-1.1);
  \node[below, font=\footnotesize] at (0,-1.2) {Administrador};

  \coordinate (profPos) at (0,-3.5);
  \draw[thick] (profPos) circle (0.2);
  \draw[thick] (0,-3.7) -- (0,-4.2);
  \draw[thick] (-0.25,-3.9) -- (0.25,-3.9);
  \draw[thick] (0,-4.2) -- (-0.2,-4.6);
  \draw[thick] (0,-4.2) -- (0.2,-4.6);
  \node[below, font=\footnotesize] at (0,-4.7) {Profesor};

  \coordinate (estPos) at (0,-7);
  \draw[thick] (estPos) circle (0.2);
  \draw[thick] (0,-7.2) -- (0,-7.7);
  \draw[thick] (-0.25,-7.4) -- (0.25,-7.4);
  \draw[thick] (0,-7.7) -- (-0.2,-8.1);
  \draw[thick] (0,-7.7) -- (0.2,-8.1);
  \node[below, font=\footnotesize] at (0,-8.2) {Estudiante};

  % ---- CASOS DE USO ----
  \node[usecase] (login) at (6,-1) {Iniciar Sesi\'on};
  \node[usecase] (genqr) at (6,-2.5) {Generar C\'odigo QR};
  \node[usecase] (scanqr) at (6,-4) {Escanear C\'odigo QR};
  \node[usecase] (gestusr) at (6,-5.5) {Gestionar Usuarios};
  \node[usecase] (config) at (6,-7) {Configurar Par\'ametros};
  \node[usecase] (report) at (6,-8.5) {Consultar Reportes};

  % ---- LÍMITES DEL SISTEMA ----
  \draw[dashed] (4,0.5) rectangle (8,-9.5);
  \node[above, font=\small\bfseries] at (6,0.5) {SISTEMA QR};

  % ---- ASOCIACIONES ----
  \draw (adminPos) -- (login);
  \draw (profPos) -- (login);
  \draw (estPos) -- (login);
  \draw (profPos) -- (genqr);
  \draw (estPos) -- (scanqr);
  \draw (adminPos) -- (gestusr);
  \draw (profPos) -- (config);
  \draw (adminPos) -- (report);

\end{tikzpicture}
\end{center}

\newpage

\section{Especificaciones de Casos de Uso}

A continuaci\'on se detallan las especificaciones de los casos de uso identificados.

% CU-01
\begin{center}
\begin{tabular}{|p{2.9cm}|p{10.4cm}|}
\hline
\multicolumn{2}{|c|}{\textbf{CU-01: Iniciar Sesi\'on}} \\
\hline
\textbf{Actores}        & Administrador, Profesor, Estudiante \\
\hline
\textbf{Descripci\'on}  & Permite a los usuarios autenticarse mediante su c\'edula o nombre de usuario y contrase\~na \\
\hline
\textbf{Precondiciones} & El usuario debe poseer una cuenta activa registrada en el sistema \\
\hline
\textbf{Postcondiciones}& El sistema redirige al usuario al panel correspondiente seg\'un su rol \\
\hline
\textbf{Flujo\newline principal}
  & 1. El actor ingresa su c\'edula y contrase\~na \newline
    2. El sistema valida que los campos no est\'en vac\'ios \newline
    3. El sistema consulta la base de datos y verifica la contrase\~na \newline
    4. El sistema verifica que la cuenta est\'e activa \newline
    5. El sistema inicia la sesi\'on y redirige seg\'un el rol \\
\hline
\textbf{Flujos\newline alternos}
  & \textbf{FA-01:} Contrase\~na incorrecta --- el sistema muestra un mensaje de error \newline
    \textbf{FA-02:} Cuenta inactiva --- el sistema deniega el acceso \newline
    \textbf{FA-03:} Contrase\~na temporal --- redirige al cambio obligatorio de contrase\~na \\
\hline
\end{tabular}
\vspace{8pt}
\end{center}

% CU-02
\begin{center}
\begin{tabular}{|p{2.9cm}|p{10.4cm}|}
\hline
\multicolumn{2}{|c|}{\textbf{CU-02: Generar C\'odigo QR}} \\
\hline
\textbf{Actores}        & Profesor \\
\hline
\textbf{Descripci\'on}  & Permite al profesor generar un c\'odigo QR \'unico y temporal para una sesi\'on de clase espec\'ifica \\
\hline
\textbf{Precondiciones} & El profesor debe haber iniciado sesi\'on y tener al menos un curso asignado en el periodo activo \\
\hline
\textbf{Postcondiciones}& Se crea un registro de sesi\'on QR con c\'odigo \'unico y fecha de expiraci\'on \\
\hline
\textbf{Flujo\newline principal}
  & 1. El profesor selecciona el curso y la semana \newline
    2. El sistema genera un c\'odigo \'unico de 32 caracteres hexadecimales criptogr\'aficos \newline
    3. El sistema calcula la fecha de expiraci\'on seg\'un la configuraci\'on del profesor \newline
    4. El sistema presenta el QR en pantalla para que el profesor lo proyecte \newline
    5. El sistema inicia la rotaci\'on autom\'atica seg\'un el intervalo configurado \\
\hline
\textbf{Flujos\newline alternos}
  & \textbf{FA-01:} Sin configuraci\'on previa --- el sistema usa valores por defecto (5\,min, 30\,s) \\
\hline
\end{tabular}
\vspace{8pt}
\end{center}

% CU-03
\begin{center}
\begin{tabular}{|p{2.9cm}|p{10.4cm}|}
\hline
\multicolumn{2}{|c|}{\textbf{CU-03: Escanear C\'odigo QR}} \\
\hline
\textbf{Actores}        & Estudiante \\
\hline
\textbf{Descripci\'on}  & Permite al estudiante registrar su asistencia escaneando el c\'odigo QR generado por el profesor \\
\hline
\textbf{Precondiciones} & El estudiante debe haber iniciado sesi\'on en su dispositivo m\'ovil y estar inscrito en el curso correspondiente \\
\hline
\textbf{Postcondiciones}& Se registra la asistencia del estudiante con estado ``presente'' para la sesi\'on y semana correspondiente \\
\hline
\textbf{Flujo\newline principal}
  & 1. El estudiante abre el m\'odulo de escaneo desde su dispositivo m\'ovil \newline
    2. El sistema activa la c\'amara del dispositivo para capturar el c\'odigo QR \newline
    3. El estudiante enfoca el QR proyectado por el profesor \newline
    4. El sistema decodifica el QR y extrae el c\'odigo \'unico de sesi\'on \newline
    5. El sistema valida que el QR exista y no haya expirado \newline
    6. El sistema registra la asistencia como ``presente'' y muestra confirmaci\'on \\
\hline
\textbf{Flujos\newline alternos}
  & \textbf{FA-01:} QR expirado --- el sistema rechaza el registro y notifica al estudiante \newline
    \textbf{FA-02:} QR inv\'alido --- el sistema muestra un mensaje de error \newline
    \textbf{FA-03:} Estudiante no inscrito --- el sistema rechaza el registro \\
\hline
\end{tabular}
\vspace{8pt}
\end{center}

% CU-04
\begin{center}
\begin{tabular}{|p{2.9cm}|p{10.4cm}|}
\hline
\multicolumn{2}{|c|}{\textbf{CU-04: Gestionar Usuarios}} \\
\hline
\textbf{Actores}        & Administrador \\
\hline
\textbf{Descripci\'on}  & Permite al administrador crear, modificar, activar o desactivar cuentas de profesores y estudiantes \\
\hline
\textbf{Precondiciones} & El administrador debe haber iniciado sesi\'on \\
\hline
\textbf{Postcondiciones}& Se crean, modifican o actualizan los registros de usuarios seg\'un la operaci\'on realizada \\
\hline
\textbf{Flujo\newline principal}
  & 1. El administrador accede al m\'odulo de gesti\'on de usuarios \newline
    2. El sistema presenta la lista de usuarios con opciones de b\'usqueda y filtros \newline
    3. El administrador selecciona la acci\'on (crear, editar, activar o desactivar) \newline
    4. Para crear, ingresa los datos del nuevo usuario (c\'edula, nombre, correo, rol, carrera) \newline
    5. El sistema valida que la c\'edula no est\'e duplicada y genera username y contrase\~na temporal \newline
    6. El sistema registra el nuevo usuario en la base de datos \\
\hline
\textbf{Flujos\newline alternos}
  & \textbf{FA-01:} C\'edula duplicada --- el sistema muestra un mensaje de error \newline
    \textbf{FA-02:} Datos inv\'alidos --- el sistema solicita la correcci\'on \\
\hline
\end{tabular}
\vspace{8pt}
\end{center}

% CU-05
\begin{center}
\begin{tabular}{|p{2.9cm}|p{10.4cm}|}
\hline
\multicolumn{2}{|c|}{\textbf{CU-05: Configurar Par\'ametros de Asistencia}} \\
\hline
\textbf{Actores}        & Profesor \\
\hline
\textbf{Descripci\'on}  & Permite al profesor ajustar los par\'ametros de expiraci\'on y rotaci\'on del QR \\
\hline
\textbf{Precondiciones} & El profesor debe haber iniciado sesi\'on \\
\hline
\textbf{Postcondiciones}& Se actualizan los valores de configuraci\'on del profesor en la base de datos \\
\hline
\textbf{Flujo\newline principal}
  & 1. El profesor accede al m\'odulo de configuraci\'on \newline
    2. El sistema muestra los valores actuales de expiraci\'on y rotaci\'on \newline
    3. El profesor modifica los par\'ametros seg\'un su preferencia \newline
    4. El sistema valida que los valores est\'en dentro de los rangos permitidos (1--60\,min, 10--300\,s) \newline
    5. El sistema guarda la configuraci\'on actualizada \\
\hline
\textbf{Flujos\newline alternos}
  & \textbf{FA-01:} Valores fuera de rango --- el sistema indica el rango permitido y solicita correcci\'on \\
\hline
\end{tabular}
\vspace{8pt}
\end{center}

% CU-06
\begin{center}
\begin{tabular}{|p{2.9cm}|p{10.4cm}|}
\hline
\multicolumn{2}{|c|}{\textbf{CU-06: Consultar Reportes de Asistencia}} \\
\hline
\textbf{Actores}        & Administrador \\
\hline
\textbf{Descripci\'on}  & Permite al administrador consultar reportes detallados de asistencia filtrados por curso, semana o estudiante \\
\hline
\textbf{Precondiciones} & El administrador debe haber iniciado sesi\'on \\
\hline
\textbf{Postcondiciones}& El sistema muestra los reportes solicitados sin realizar modificaciones en los datos \\
\hline
\textbf{Flujo\newline principal}
  & 1. El administrador accede al m\'odulo de reportes \newline
    2. El sistema presenta opciones de filtrado: curso, semana, periodo y estudiante \newline
    3. El administrador selecciona los filtros deseados \newline
    4. El sistema consulta la base de datos y genera una tabla con resultados \newline
    5. El sistema muestra el reporte con asistencias, inasistencias y porcentaje de ausentismo \\
\hline
\textbf{Flujos\newline alternos}
  & \textbf{FA-01:} Sin resultados --- el sistema indica que no hay datos disponibles \\
\hline
\end{tabular}
\vspace{8pt}
\end{center}

\newpage

% ============================================================
% REFERENCIAS
% ============================================================
\chapter*{REFERENCIAS}
\addcontentsline{toc}{chapter}{REFERENCIAS}
\markboth{REFERENCIAS}{}

\begin{enumerate}
\item Booch, G., Rumbaugh, J., \& Jacobson, I. (2005). \textit{El lenguaje unificado de modelado: Gu\'ia del usuario} (2.a ed.). Addison-Wesley.
\item Object Management Group. (2017). \textit{Unified Modeling Language (UML), version 2.5.1}. OMG.
\item Pressman, R. S. (2014). \textit{Ingenier\'ia del software: Un enfoque pr\'actico} (7.a ed.). McGraw-Hill.
\item Sommerville, I. (2011). \textit{Ingenier\'ia de software} (9.a ed.). Pearson.
\end{enumerate}

\end{document}
